
\title{Concise axiomatizations of\\ intuitionistic primitive recursive arithmetic}
\author{}
\date{November 2026}
\maketitle

\begin{abstract}
Intuitionistic first-order primitive recursive arithmetic is bi-interpretable
with a calculus that has variables, zero, successor and one iterator as its
terms, equations and two binders as its formulas, and eight inference rules.
Both numerical domains are whole, equality is actual equality, and both
composite maps have graph $x=y$. This article states that calculus and the
bridge once, and then collects every refinement of it that the research
packages of the \texttt{ConciseIPRA} directory of the \texttt{ProveIt}
repository contain: changes of the rule basis, guarded formula languages,
rules that merge computation with equality, normal forms, clause calculi,
rules that discharge a witness, introductions that take no witness, other
recursors, obstruction results, and proof-certificate formats with their
checkers. Each definition, rule inventory, theorem and proof is given once, in
one notation. Statements that the packages leave unproved are collected at the
end of each chapter.
\end{abstract}

\tableofcontents

\chapter*{Introduction}
\addcontentsline{toc}{chapter}{Introduction}

\section*{The question}

The question behind all of this work is how small a formal system can be
while it still is, in the strongest reasonable sense, intuitionistic primitive
recursive arithmetic $\iPRA$. ``Is'' means a bi-interpretation in which the
numerical domain of each interpretation is the whole domain, equality is
interpreted by equality, and each composite of the two numerical maps is the
identity, provably. ``Small'' has two readings. The mathematical one counts
term constructors, formula constructors and inference rules, and asks which of
them can be removed, restricted or fused. The empirical one fixes a
measurement protocol for complete proof checkers and compares their sizes
(\cref{base:sec:measure}). The two readings do not always agree, and the text
keeps them apart: a smaller rule basis can need a larger checker, and several
of the smallest checkers implement an unchanged rule basis.

\section*{How the article is organized}

\Cref{base:sec:terms} to \cref{base:sec:measure} of the first chapter state
the common core: the terms, the native formulas and their ordinary reading,
the eight rules of the native scalar core $\Nat$, its structural properties
and derived connectives, the ordinary theory $\Teq$, the arithmetic inside
$\Teq$, and the bridge to $\iPRA$. No other chapter restates this material.

The next six chapters stay on iterator terms. The second treats refinements
of the eight rules on the unchanged language: the introductions, the
transport rule, uniqueness of iteration in place of induction, and the
progressive reading of the closure binder. The third treats guards: formula
languages in which the shape of a matrix determines the mode of its
introduction. The fourth treats rules that merge computation with equality and
calculi on normal forms. The fifth treats clause calculi, in which contexts
hold reusable clauses and the closure binder is not a formula. The sixth
treats rules that fuse equality reasoning with the discharge of an opened
witness, and restricted substitution rules. The seventh treats introductions
that take no witness.

Two chapters change the term language: recursors with a built-in seed, and
canonical and counter-aware recursors in which the successor is defined.
The tenth chapter collects the obstruction and independence results, with the
models and invariants that prove them. The last chapter treats
proof-certificate formats and checkers, and records the measured sizes.

Each chapter ends with a table of the checker sizes that its packages report
and with a section of further questions: claims stated without proof,
conjectures, and what was tried and not achieved.

\section*{Conventions}

The notation of the first chapter is used throughout. Names of calculi other
than $\Nat$, $\Teq$ and $\iPRA$ are local to the chapter that introduces
them; the same letter can denote different calculi in different chapters.
Every chapter compares its calculi with the eight-rule core of the first
chapter. A research package is cited by its directory name, as in
\pkg{witness\_rule}; \cref{prov:sec:packages} lists every package with the
chapters that draw on it. Where a result is formalized in Lean in the
repository, the name of the Lean declaration is given in typewriter type; the
status of the formalization is recorded in the README of the directory and is
not repeated here in full. Editorial repairs of the source reports are stated
where they occur.
